\documentclass[12pt,a4paper]{article}
\usepackage{expediente}
\begin{document}
% =============================================================================
% FICHA DE DATOS DEL EXPEDIENTE --- actualizado 20 ago 2026
% =============================================================================

% --- VENDEDOR / TITULAR REGISTRAL -------------------------------------------
\newcommand{\VendedorNombres}{RICHAR DARWIN}
\newcommand{\VendedorApellidos}{CUTIRE CONDORI}
\newcommand{\VendedorNombreCompleto}{RICHAR DARWIN CUTIRE CONDORI}
\newcommand{\VendedorDNI}{42610690}
\newcommand{\VendedorEstadoCivil}{soltero, sin unión de hecho vigente}
\newcommand{\VendedorOcupacion}{\faltante{[ocupación de Richar]}}
\newcommand{\VendedorDomicilio}{\faltante{[domicilio actual de Richar]}}
\newcommand{\VendedorTelefono}{\faltante{[teléfono de Richar]}}
\newcommand{\VendedorCorreo}{\faltante{[correo]}}
\newcommand{\VendedorNacionalidad}{peruano}
\newcommand{\VendedorFechaNac}{\faltante{[fecha de nacimiento]}}
\newcommand{\VendedorDNIEmision}{\faltante{[fecha de emisión del DNI de Richar --- la pide SUNARP]}}

% --- CONVIVIENTE: NO CORRESPONDE -------------------------------------------
\newcommand{\ConvivienteNombre}{NO CORRESPONDE}
\newcommand{\ConvivienteDNI}{NO CORRESPONDE}
\newcommand{\ConvivienteDesde}{NO CORRESPONDE}
\newcommand{\ConvivienteInscrita}{no existe unión de hecho vigente}
\newcommand{\ConvivienteDomicilio}{NO CORRESPONDE}

% --- COMPRADORES (copropiedad en partes iguales) ---------------------------
\newcommand{\CompradorUnoNombre}{NICO ALVARO CUTIRE ARCE}
\newcommand{\CompradorUnoDNI}{48325017}
\newcommand{\CompradorUnoEmision}{01 de enero de 2024}
\newcommand{\CompradorUnoEstadoCivil}{soltero}
\newcommand{\CompradorUnoOcupacion}{ingeniero informático}
\newcommand{\CompradorUnoDomicilio}{Jr.\ Wiracocha, Sicuani, provincia de Canchis, Cusco, Perú}
\newcommand{\CompradorUnoTelefono}{984715108}

\newcommand{\CompradorDosNombre}{LUCY CUTIRE ARCE}
\newcommand{\CompradorDosDNI}{47478852}
\newcommand{\CompradorDosEmision}{09 de enero de 2026}
\newcommand{\CompradorDosEstadoCivil}{soltera}
\newcommand{\CompradorDosOcupacion}{contador}
\newcommand{\CompradorDosDomicilio}{Comunidad Hanocca, distrito de Layo, provincia de Canas, Cusco, Perú}
\newcommand{\CompradorDosTelefono}{953440875}

\newcommand{\CompradorNombres}{NICO ALVARO y LUCY}
\newcommand{\CompradorApellidos}{CUTIRE ARCE}
\newcommand{\CompradorNombreCompleto}{NICO ALVARO CUTIRE ARCE y LUCY CUTIRE ARCE}
\newcommand{\CompradorDNI}{48325017 y 47478852}
\newcommand{\CompradorEstadoCivil}{solteros}
\newcommand{\CompradorConyuge}{NO CORRESPONDE}
\newcommand{\CompradorConyugeDNI}{NO CORRESPONDE}
\newcommand{\CompradorOcupacion}{ingeniero informático y contador, respectivamente}
\newcommand{\CompradorDomicilio}{Sicuani (Canchis) y Layo (Canas), Cusco}
\newcommand{\CompradorTelefono}{984715108 / 953440875}
\newcommand{\CompradorCorreo}{\faltante{[correo de contacto]}}
\newcommand{\CompradorNacionalidad}{peruanos}
\newcommand{\PorcentajeCopropiedad}{50\,\% cada uno}

% --- INTERMEDIARIO / OCUPANTE DEL REMANENTE --------------------------------
\newcommand{\IntermediarioNombre}{LUCIO GIL CUTIRE PARI}
\newcommand{\IntermediarioDNI}{42467850}
\newcommand{\IntermediarioEstadoCivil}{\faltante{[estado civil de Lucio]}}
\newcommand{\IntermediarioOcupacion}{\faltante{[ocupación]}}
\newcommand{\IntermediarioDomicilio}{\faltante{[domicilio de Lucio]}}
\newcommand{\IntermediarioTelefono}{\faltante{[teléfono de Lucio]}}
\newcommand{\IntermediarioParentesco}{familiar (apellido Cutire)}

% --- PREDIO MATRIZ ----------------------------------------------------------
\newcommand{\PartidaMatriz}{\faltante{[N.º de partida electrónica SUNARP --- pendiente]}}
\newcommand{\OficinaRegistral}{Oficina Registral de Puerto Maldonado}
\newcommand{\AreaMatriz}{\faltante{[área total inscrita; las partes refieren del orden de 100 ha]}}
\newcommand{\UnidadCatastral}{\faltante{[código catastral / UC, si existe]}}
\newcommand{\NombreFundo}{\faltante{[nombre del fundo, si tiene]}}
\newcommand{\KmCarretera}{Interoceánica Sur, al este del Puente Jayave (el puente queda al oeste del predio, sobre el río; no debajo del lote)}

% --- PREDIO A INDEPENDIZAR --------------------------------------------------
\newcommand{\AreaIndependizarHa}{10{,}00}
\newcommand{\AreaIndependizarMetros}{100 000{,}00}
\newcommand{\PerimetroM}{2 200{,}58}
\newcommand{\FrenteM}{100{,}29}
\newcommand{\FondoM}{1 000{,}00}

\newcommand{\PUnoLat}{$-12{,}912452$}
\newcommand{\PUnoLon}{$-70{,}170058$}
\newcommand{\PDosLat}{$-12{,}912497$}
\newcommand{\PDosLon}{$-70{,}170981$}
\newcommand{\PTresLat}{$-12{,}903469$}
\newcommand{\PTresLon}{$-70{,}171438$}
\newcommand{\PCuatroLat}{$-12{,}903424$}
\newcommand{\PCuatroLon}{$-70{,}170515$}

\newcommand{\PUnoE}{373 065{,}00}
\newcommand{\PUnoN}{8 572 256{,}00}
\newcommand{\PDosE}{372 965{,}00}
\newcommand{\PDosN}{8 572 251{,}00}
\newcommand{\PTresE}{372 911{,}00}
\newcommand{\PTresN}{8 573 249{,}00}
\newcommand{\PCuatroE}{373 011{,}00}
\newcommand{\PCuatroN}{8 573 255{,}00}

\newcommand{\FechaPago}{15 de agosto de 2018}
\newcommand{\MontoPago}{S/\ 20 000{,}00}
\newcommand{\MontoPagoLetras}{VEINTE MIL Y 00/100 SOLES}
\newcommand{\FormaPagoHistorica}{dinero en efectivo, sin medio de pago del sistema financiero}

% Nombres genéricos --- sustituir por los reales al ir a notaría
\newcommand{\NotarioNombre}{Dr.\ JUAN CARLOS RÍOS SALINAS (nombre genérico a sustituir)}
\newcommand{\NotariaCiudad}{Puerto Maldonado}
\newcommand{\ColegiaturaAbogado}{CAL Madre de Dios N.º 0000 (genérico)}
\newcommand{\AbogadoNombre}{Abog.\ ANA MARÍA QUISPE HUAMÁN (nombre genérico a sustituir)}
\newcommand{\IngenieroNombre}{Ing.\ PEDRO ANTONIO MAMANI CONDORI (nombre genérico a sustituir)}
\newcommand{\IngenieroCIP}{CIP 000000 (genérico)}
\newcommand{\Municipalidad}{Municipalidad Distrital de Inambari}
\newcommand{\Provincia}{Tambopata}
\newcommand{\Departamento}{Madre de Dios}
\newcommand{\Distrito}{Inambari}

\newcommand{\LugarOtorgamiento}{Puerto Maldonado}
\newcommand{\DiaOtorgamiento}{\faltante{[día]}}
\newcommand{\MesOtorgamiento}{\faltante{[mes]}}
\newcommand{\AnioOtorgamiento}{2026}

\newcommand{\TestigoUno}{\faltante{[testigo 1: nombres, DNI, domicilio]}}
\newcommand{\TestigoDos}{\faltante{[testigo 2: nombres, DNI, domicilio]}}

\InicioDocumento{Solicitud de liquidación y constancia del Impuesto de Alcabala}{DOC-07}

\noindent \textbf{SEÑOR ALCALDE / GERENCIA DE RENTAS}\\
\textbf{\MakeUppercase{\Municipalidad}} \quad (o SAT / órgano que administre el tributo en el distrito de \Distrito)

\vspace{0.3em}
\noindent \CompradorNombreCompleto, DNI \CompradorDNI, en calidad de \textbf{sujeto pasivo del Impuesto de Alcabala} por la adquisición que se formalizará, con domicilio en \CompradorDomicilio, atentamente \textbf{digo}:

\Clausula{Petitorio}
Solicito la \textbf{liquidación del Impuesto de Alcabala} y la \textbf{constancia de pago} (o de inafectación, si el tramo no gravado cubre la base imponible), para su presentación ante notario y SUNARP, respecto de la transferencia de un predio rústico de \textbf{10,00 hectáreas} que independizará y venderá el señor \VendedorNombreCompleto.

\Clausula{Datos de la transferencia}
\renewcommand{\arraystretch}{1.28}
\noindent\begin{tabularx}{\textwidth}{>{\bfseries}p{5.3cm}X}
Transferente & \VendedorNombreCompleto, DNI \VendedorDNI \\
Adquirente & \CompradorNombreCompleto, DNI \CompradorDNI \\
Acto & Compraventa (título oneroso), no donación \\
Fecha histórica del trato & \FechaPago\ (documento privado); fecha de formalización: \DiaOtorgamiento\ de \MesOtorgamiento\ de \AnioOtorgamiento \\
Precio convenido & \MontoPago\ (\MontoPagoLetras) \\
Valor de autoavalúo & \faltante{[valor de autoavalúo ajustado que determine la Municipalidad]} \\
Base imponible esperada & El \textbf{mayor} entre el valor de transferencia y el autoavalúo, conforme al art.\ 24 de la Ley de Tributación Municipal \\
Deducción & Primeras 10 UIT del ejercicio de la transferencia \\
Inmueble & 10,00 ha a independizar del predio de partida \PartidaMatriz, \Distrito, \Provincia, \Departamento \\
\end{tabularx}

\Clausula{Fundamento}
Arts.\ 21 a 26 del TUO de la Ley de Tributación Municipal. El alcabala grava transferencias de predios urbanos o rústicos a título oneroso o gratuito. Lo paga el comprador. El notario y el registrador requieren acreditar el pago (o la inafectación) para autorizar e inscribir.

\Clausula{Declaración}
El adquirente declara que el precio histórico fue \MontoPago\ pagado en efectivo en 2018 y que \textbf{no solicita} se liquide el acto como donación. Si la Municipalidad toma un autoavalúo superior al precio, acepta que la base sea ese mayor valor.

\vspace{0.35em}
\noindent POR TANTO:\\
Pido la liquidación, el comprobante y la constancia para notaría.

\LugarFecha

\noindent
\BloqueFirma{EL ADQUIRENTE (ALCABALA)}{\CompradorNombreCompleto}{\CompradorDNI}
\hfill
\BloqueFirmaSimple{ENTERADO EL TRANSFERENTE}{\VendedorNombreCompleto\\DNI \VendedorDNI}

\vspace{0.9em}
\noindent{\small\textit{Uso municipal --- liquidación}}\\
{\small
\begin{tabularx}{\textwidth}{|p{5.6cm}|X|}
\hline
Valor de transferencia declarado & \MontoPago \\
\hline
Autoavalúo ajustado & \linea[6cm] \\
\hline
Base imponible (el mayor) & \linea[6cm] \\
\hline
(-) 10 UIT del ejercicio & \linea[6cm] \\
\hline
Base neta gravada & \linea[6cm] \\
\hline
Impuesto 3\,\% & \linea[6cm] \\
\hline
Inafectación / monto a pagar & \linea[6cm] \\
\hline
N.º liquidación / constancia & \linea[6cm] \\
\hline
Fecha de pago & \linea[6cm] \\
\hline
\end{tabularx}}

\end{document}
